\section{Method}

\subsection{Overview}

Let $I$ denote a query image and let $B(I)$ be the feature map produced by a relocalization backbone. Standard SCR methods learn a scene-specific function that maps image features to 3D scene coordinates and then recover pose through a geometric solver. \shortmethod{} augments this pipeline with a scene memory $M=\{(p_i, f_i)\}_{i=1}^{N}$, where $p_i$ is a 3D memory point and $f_i$ is its associated feature. The memory may come from MapAnything features, ACE-FCN feature-space extraction, or GLACE-compatible representations. A compressor $C_\theta$ maps this memory into $K$ latent tokens,
\begin{equation}
    Z = C_\theta(M), \qquad Z \in \mathbb{R}^{K \times d},
\end{equation}
where $K$ is much smaller than the number of memory points. A fusion module $F_\phi$ then conditions query features on the compressed memory,
\begin{equation}
    \tilde{B}(I) = F_\phi(B(I), Z),
\end{equation}
and the coordinate regressor predicts scene coordinates from $\tilde{B}(I)$.

\subsection{Memory Compression}

The current implementation uses a geometric latent memory compressor, GeoLMC, which attends from learnable latent tokens to pooled scene memory features. The memory loader supports both pooled format and BSE-style memory files and normalizes scene points around a scene center. The compressed tokens are shared across query images of the same scene during evaluation, so memory compression is performed once per scene rather than once per frame.

\subsection{One-Level Two-Stage Training}

The one-level configuration is used for the main Indoor6 DINOv2+MapAnything result. The backbone is DINOv2, and the memory source is MapAnything-derived scene memory. Training alternates between two stages. The first stage aligns the compressor and fusion module with query features, while the second stage refines the coordinate regressor using reprojection loss. This setting corresponds to the ACE-G-style LMC flow in the codebase and uses raw backbone buffers so that memory fusion remains part of the trainable path.

\subsection{Two-Level Two-Stage Training}

The two-level configuration separates local memory learning from global-local regression. In the first level, ACE-FCN feature-space memory is extracted from a frozen ACE encoder. This memory has stride 8 and 512-dimensional features, and its memory points are derived from patch-level scene coordinates. A local ACE-FCN-LMC checkpoint is trained without GLACE global conditioning. In the second level, this local stack is loaded into a GLACE-style regressor and combined with global features through concatenation, residual, or FiLM-style conditioning. This design tests whether a local latent memory representation can serve as a plug-in module for a different global-local relocalization architecture.

\subsection{Visibility-Guided Configuration Selection}

Global features are useful only when the memory and query distribution support reliable global conditioning. We therefore define \gvcs{}, a visibility-guided configuration selection rule. The implementation estimates the front-facing visibility of memory points and computes a routing score based on the median front ratio minus the fraction of low-pose front visibility. When the requested configuration is global but the score falls below a threshold, the system falls back to local or hierarchical conditioning. \gvcs{} is intentionally simple and deterministic in the current version; it should be interpreted as a reliability-guided routing rule rather than a learned scene-understanding policy.

\begin{algorithm}[t]
\caption{Visibility-guided configuration selection}
\label{alg:gvcs}
\begin{algorithmic}[1]
\STATE Estimate memory front visibility for the current scene.
\STATE Compute $s = \mathrm{median\_front\_ratio} - \mathrm{low\_pose\_front\_ratio\_fraction}$.
\IF{requested mode is global and $s < \tau$}
    \STATE Select local or hierarchical fallback mode.
\ELSE
    \STATE Keep requested mode.
\ENDIF
\STATE Train or evaluate with the selected conditioning configuration.
\end{algorithmic}
\end{algorithm}
